\centerline{\bf \Huge Комби разнобой}

\begin{flushright}
        \scriptsize{}
\end{flushright}

\problem{1}
В олимпиаде принимало участие семеро школьников, каждый из которых получил от 0 до 11 баллов. Докажите, что можно выбрать два непересекающихся множества школьников, в которых поровну ребят, и в которых суммы набранных баллов равны.

\problem{2}
Изначально по кругу расставлены 40 синих, 30 красных и 20 зеленых фишек, причем фишки каждого цвета идут подряд. За ход можно поменять местами стоящие рядом синюю и красную фишки, или стоящие рядом синюю и зеленую фишки. Можно ли за несколько таких операций добиться того, чтобы любые две стоящие рядом фишки были разных цветов?

\problem{3}
Двое игроков по очереди расставляют в каждой из 24 клеток поверхности куба $2 \times 2 \times 2$ числа $1,2,3, \ldots, 24$ (каждое число можно ставить один раз). Второй игрок хочет, чтобы суммы чисел в клетках каждого кольца из 8 клеток, опоясывающего куб, были одинаковыми. Сможет ли первый игрок ему помешать?

\problem{4}
В компании 100 человек. Оказалось, что любых 98 из них можно разбить на 49 пар знакомых. Какое наименьшее число пар знакомых может быть в этой компании?

\problem{5}
На окружности отмечено 150 серых, 151 бурая и 152 малиновых точки таким образом, что никакие две одноцветные точки не стоят рядом. Докажите, что найдётся бурая точка, у которой оба соседа - малиновые.

\problem{6}
На экзамен пришли 100 студентов. Преподаватель по очереди задаёт каждому студенту один вопрос: «Сколько из 100 студентов получат оценку «сдал» к концу экзамена?». В ответ студент называет целое число. Сразу после получения ответа преподаватель объявляет всем, какую оценку получил студент: «сдал» или «не сдал».

После того, как все студенты получат оценку, придет инспектор и проверит, есть ли студенты, которые дали правильный ответ, но получили оценку «не сдал». Если хотя бы один такой студент найдётся, то преподаватель будет отстранен от работы, а оценки всех студентов заменят на «сдал». В противном случае никаких изменений не произойдёт.

(a) Придумайте стратегию, которая гарантирует всем студентам оценку «сдал».

(b) Докажите, что эта стратегия единственная.